% Generated from locale/tr/content/history/biographies/gerhard-gentzen.tex; do not hand-edit.


\olfileid[tr]{his}{bio}{gen}

\olsection{Gerhard Gentzen}

\olphoto{gentzen-gerhard}{Gerhard Gentzen}

Gerhard Gentzen öncelikle yapısal ispat teorisinin kurucusu, özellikle doğal
tümdengelim ve sekant hesabı !!{derivation} sistemlerini geliştiren kişi
olarak tanınır. 24 Kasım 1909'da Almanya'nın Greifswald kentinde doğdu.
Gerhard hazırlık okuluna başlamadan önce üç yıl evde eğitim gördü; okula
başladığında eğitim bakımından sınıf arkadaşlarının çoğunun gerisindeydi.
Buna rağmen parlak bir öğrenciydi ve matematikte güçlü bir yetenek gösterdi.
İlgi alanları çeşitliydi; örneğin annesi için şiirler ve okul tiyatrosu için
oyunlar da yazardı.

Gentzen üniversite eğitimine Greifswald Üniversitesi'nde başladı, daha sonra
Göttingen, Münih ve Berlin'e geçti. Doktorasını 1933'te Göttingen
Üniversitesi'nde Hermann Weyl yönetiminde aldı. (Çalışmasının büyük bölümünü
Paul Bernays yönetmişti; fakat Naziler onu üniversiteden uzaklaştırdı.)
Gentzen 1934'te David Hilbert'in asistanı olarak çalışmaya başladı. Aynı yıl
\emph{Untersuchungen über das logische Schließen I--II [Mantıksal
Tümdengelim Üzerine İncelemeler I--II]} başlıklı makalelerinde sekant hesabı
ve doğal tümdengelim !!{derivation} sistemlerini geliştirdi. 1936'da Peano
aksiyomlarının tutarlılığını ispatladı.

Gentzen'in Nazilerle ilişkisi karmaşıktır. Danışmanı Bernays Almanya'yı terk
etmek zorunda bırakıldığı sırada Gentzen, Nazi paramiliter örgütü SA'nın
üniversite koluna katıldı. Birçok Alman gibi Nazi Partisi'ne üyeydi. Savaş
sırasında hava istihbarat biriminde haberleşme subayı olarak görev yaptı.
Ancak 1942'de sinir krizi nedeniyle görevden alındı. Gentzen'in gerçekten
Nazi Partisi'ne bağlı olup olmadığı ya da partiye akademik başarısını güvence
altına almak için katılıp katılmadığı belirsizdir.

Gentzen 1943'te Prag Alman Üniversitesi Matematik Enstitüsü'nden gelen
akademik görev teklifini kabul etti. Ancak 1945'te Prag halkı Alman işgaline
karşı ayaklandı. Sovyet kuvvetleri kente girip üniversitedeki bütün
profesörleri tutukladı. Nazi örgütlerine üyeliği nedeniyle Gentzen zorunlu
çalışma kampına götürüldü. 4 Ağustos 1945'te, 35 yaşındayken hücresinde
yetersiz beslenmeden öldü.

\begin{reading}
Gentzen'in kapsamlı bir yaşam öyküsü için \citet{Menzler-Trott2007}
kaynağına bakın. Nazi yönetimi altındaki matematikçileri konu alan ve
Gentzen'in yaşamına kısaca değinen ilginç bir çalışma
\citet{Segal2014}'tür. Gentzen'in mantıksal tümdengelim makaleleri özgün
Almancalarıyla mevcuttur \citep{Gentzen1935a,Gentzen1935b}.
\citet{Gentzen1969}, Gentzen'in makalelerinin İngilizce çevirilerini,
biyografik bir taslakla birlikte tek bir ciltte toplamıştır.
\end{reading}

